\chapter{Historique des versions}
Cette annexe résume les grandes étapes du développement de \soft, des plus récentes aux plus anciennes. Les
numéros sont ceux des versions de développement remises pour essai.

\section*{Version 40}
\begin{itemize}
  \item Interface entièrement en anglais.
  \item SID : la colonne Mod revient à la seule voix source (« Sync V1 », « Sync+Ring V1 »).
\end{itemize}
\section*{Version 39}
\begin{itemize}
  \item Vue PC Engine : forme de l'onde décrite en mots (sinus, triangle, scie montante ou descendante, carré,
        impulsion), numéro de table et amplitude.
\end{itemize}
\section*{Version 38}
\begin{itemize}
  \item Vue PC Engine : barre de niveau et barres de panoramique.
  \item SID : la voie des samples se place sous les voix de sa puce ; les trois voix ne contiennent plus les
        échantillons joués par le volume.
  \item AdLib OPL/OPL2 : neuf voix en 3~$\times$~3, voies de batterie masquées par défaut.
\end{itemize}
\section*{Version 37}
\begin{itemize}
  \item PC Engine : sortie par voix et état du HuC6280 dans le cœur Ootake ; vue dédiée avec la table d'ondes de
        chaque voix et tous ses registres.
\end{itemize}
\section*{Version 36}
\begin{itemize}
  \item Plein écran portrait : encoche du haut et coins du bas évités, tableau de bord correctement placé.
  \item SID : la colonne Mod indique la voix source de la synchronisation et de l'anneau.
\end{itemize}
\section*{Version 35}
\begin{itemize}
  \item Réglages de hauteur de l'oscilloscope et de la réponse du filtre.
  \item SID : voie des samples (\$D418) en option.
  \item Panneau en direct : page SNES.
\end{itemize}
\section*{Versions 33 et 34}
\begin{itemize}
  \item Listes de lecture multiples, favoris, ajout depuis le lecteur, le navigateur et les listes.
  \item File d'attente « Jouer ensuite », import et export M3U, bouton $\equiv$+ dans le lecteur.
  \item Fichiers multi-morceaux lus directement depuis le navigateur.
  \item SID : sortie mono en option pour les morceaux 2SID et 3SID.
\end{itemize}
\section*{Versions 30 à 32}
\begin{itemize}
  \item Super Nintendo : émulation ares (au cycle près) par défaut, snes\_spc en option, pour les SPC et les SNSF.
\end{itemize}
\section*{Versions antérieures}
\begin{itemize}
  \item Vue de la réponse du filtre SID dans l'oscilloscope, sans texte de valeurs.
  \item DSP d'effets du SCSP activé : réverbération des bornes Sega et de la Saturn.
  \item Valeurs du tableau de bord en temps réel, image par image.
  \item Informations techniques du morceau (codec, fréquence, résolution, débit), bas d'écran compact.
  \item Cœurs FM au choix par famille, préréglages de plateforme et étage de sortie analogique.
  \item Vue des puces PCM YMW258 et SCSP.
  \item WAV étendu (PCM et virgule flottante).
  \item Identité LHUNA : barre de navigation, mini-lecteur, réglages en cartes, À propos.
  \item Oscilloscope Vulkan et OpenGL ES à traits fins, fenêtre adaptée à la période de chaque voix.
  \item Inspecteur FM (OPN, OPM, OPL) et vue S-DSP.
\end{itemize}
